\chapter{Koncepcja Hardware-in-the-Loop (HIL) i autorski system badawczy}

Weryfikacja niezawodności i wydajności systemów wbudowanych operujących w czasie rzeczywistym wymaga zastosowania odpowiednio zaprojektowanych środowisk testowych.
Tradycyjne metody oceny oprogramowania, opierające się wyłącznie na analizie statycznej kodu lub symulacjach programowych, nie pozwalają w pełni uwzględnić interakcji zachodzących pomiędzy rzeczywistym procesorem mikrokontrolera, jego fizycznymi peryferiami, stosem komunikacyjnym oraz asynchronicznymi zdarzeniami sieciowymi.
Z tego względu do ewaluacji stworzonego oprogramowania sieciowego wybrano metodologię \textit{Hardware-in-the-Loop} (HIL).

Realizacja stanowiska badawczego wymagała zaprojektowania w pełni zautomatyzowanego środowiska testowego, w którym komputer PC pełni funkcję jednostki nadrzędnej, natomiast mikrokontroler ESP32 stanowi urządzenie badane (\textit{ang. DUT - Device Under Test}).
Oprogramowanie sterujące, zaimplementowane w języku Python, pełni funkcję orkiestratora eksperymentów, generatora ruchu sieciowego oraz menedżera urządzenia, a także realizuje proces akwizycji i analizy wyników.
Pozwala ono na automatyczne wykonywanie zdefiniowanej macierzy testowej oraz przetwarzanie zebranych danych na użyteczne metryki jakości transmisji.

W niniejszym rozdziale przedstawiono teoretyczne założenia techniki HIL w kontekście systemów wbudowanych, a następnie zaprezentowano architekturę autorskiego stanowiska badawczego.
Omówiono również zastosowane mechanizmy automatyzacji eksperymentów, sposób generowania macierzy testowej oraz ograniczenia związane z generowaniem kontrolowanego ruchu sieciowego z wykorzystaniem komputera wyposażonego w system operacyjny ogólnego przeznaczenia (GPOS).

\section{Micro-HIL jako metoda ewaluacji sprzętu}

Technika \textit{Hardware-in-the-Loop} (HIL), określana jako testowanie \textit{sprzętu w pętli}, jest metodą testowania i walidacji systemów, w której rzeczywiste urządzenie badane współpracuje z komputerowo realizowanym modelem lub środowiskiem testowym.
W klasycznym stanowisku HIL system nadrzędny odwzorowuje w czasie rzeczywistym zachowanie środowiska, z którym urządzenie współpracuje w warunkach docelowych, umożliwiając badanie rzeczywistego sprzętu i oprogramowania w kontrolowanych oraz powtarzalnych warunkach.

Z perspektywy testowanego mikrokontrolera środowisko HIL dostarcza rzeczywistych bodźców wejściowych za pośrednictwem wykorzystywanych w systemie interfejsów komunikacyjnych.
Dzięki temu testowany układ wykonuje rzeczywisty kod aplikacyjny na docelowym sprzęcie, z wykorzystaniem jego procesora, pamięci, fizycznych peryferiów oraz rzeczywistego stosu komunikacyjnego.

Zastosowana w pracy architektura stanowi uproszczoną odmianę środowiska Hardware-in-the-Loop, którą można określić mianem Micro-HIL (MHIL).
W klasycznym środowisku HIL rzeczywisty sterownik jest integrowany z systemem czasu rzeczywistego wykonującym model symulowanego środowiska (\textit{ang .plant}), a wymiana danych odbywa się za pośrednictwem rzeczywistych interfejsów wejścia-wyjścia \cite{mathworks_hil}.
W projektowanym systemie nie jest konieczne odwzorowanie złożonego modelu fizycznego badanego środowiska.
Komputer PC pełni przede wszystkim funkcję orkiestratora eksperymentu, generatora kontrolowanego ruchu sieciowego oraz systemu akwizycji i analizy wyników, natomiast rzeczywisty mikrokontroler ESP32 stanowi urządzenie badane (DUT).

Takie ograniczenie zakresu symulowanego środowiska pozwala uprościć architekturę stanowiska, zachowując jednocześnie istotną cechę podejścia HIL, jaką jest wykonywanie testowanego oprogramowania na rzeczywistym sprzęcie.
Koncepcja uproszczonych systemów tego typu znajduje odzwierciedlenie w literaturze dotyczącej Micro Hardware-in-the-Loop, w której przedstawiono możliwość redukcji klasycznego stanowiska HIL do mniejszego i bardziej ekonomicznego systemu testowego \cite{palladino2009micro}.

\subsection{Zalety testowania zjawisk sieciowych w pętli sprzętowej}

W kontekście badań wydajności transmisyjnej, obejmujących między innymi pomiar \textit{jitteru}, przepustowości oraz strat pakietów (\textit{PDR}), zastosowanie rzeczywistego urządzenia docelowego pozwala uwzględnić czynniki, których nie można bezpośrednio zaobserwować w przypadku czystej symulacji programowej.

\begin{itemize}

    \item \textbf{Walidacja rzeczywistego stosu TCP/IP:}
          w przeciwieństwie do symulacji programowej, w której zachowanie urządzenia oraz jego stosu komunikacyjnego jest reprezentowane przez model, zastosowanie HIL umożliwia wykonanie rzeczywistego oprogramowania na docelowym mikrokontrolerze.
          Dzięki temu pomiar obejmuje również ograniczenia wynikające z rzeczywistego procesora, pamięci, peryferiów oraz implementacji stosu komunikacyjnego.

    \item \textbf{Identyfikacja ograniczeń sprzętowych:}
          w testach obejmujących krótkotrwałe zwiększenie intensywności ruchu (\textit{ang. Burst tests}) badaniu podlega nie tylko zachowanie protokołów komunikacyjnych, lecz również wpływ ograniczeń sprzętowych badanego urządzenia.
          Dotyczy to między innymi wydajności obsługi interfejsów komunikacyjnych, dostępu do pamięci oraz zasobów procesora.
          W przypadku wykorzystania dodatkowych komponentów sprzętowych analiza może również obejmować ich wpływ na zdolność mikrokontrolera do obsługi zwiększonego obciążenia.

    \item \textbf{Automatyzacja i powtarzalność:}
          komputer sterujący kontroluje parametry macierzy testowej, w tym częstotliwość generowania ruchu, wykorzystywany protokół oraz rozmiar ładunku użytecznego.
          Automatyzacja procesu pozwala ograniczyć wpływ operatora na przebieg eksperymentu oraz zwiększyć powtarzalność warunków badawczych.

\end{itemize}

\subsection{Wyzwania i ograniczenia systemów HIL}

Istotnym wyzwaniem w przypadku zastosowanego stanowiska jest fakt, że komputer PC pełniący funkcję jednostki nadrzędnej pracuje pod kontrolą systemu operacyjnego ogólnego przeznaczenia (GPOS), takiego jak Windows lub Linux.
W przeciwieństwie do systemów czasu rzeczywistego systemy te nie gwarantują deterministycznego czasu wykonania operacji.
Na moment rozpoczęcia wykonywania poszczególnych operacji wpływają między innymi mechanizmy szeregowania zadań przez system operacyjny, obciążenie procesora oraz działanie innych procesów.

Ograniczenie to ma szczególne znaczenie podczas generowania ruchu sieciowego z określonymi odstępami czasowymi.
Wykorzystanie standardowych mechanizmów usypiania wątku, takich jak \texttt{time.sleep()} w języku Python, nie gwarantuje dokładnego momentu jego wznowienia, ponieważ rzeczywisty czas oczekiwania zależy od mechanizmów planowania zadań systemu operacyjnego.
Z tego względu w krytycznych fragmentach generatora ruchu zastosowano aktywne oczekiwanie (\textit{ang. busy-wait}), którego zadaniem jest ograniczenie zmienności momentu rozpoczęcia kolejnych operacji.

Zastosowanie aktywnego oczekiwania nie zapewnia jednak deterministycznego czasu wykonania charakterystycznego dla systemów czasu rzeczywistego.
Ostateczna precyzja generowania ruchu zależy również od obciążenia systemu operacyjnego, działania stosu sieciowego, wykorzystywanego interfejsu komunikacyjnego oraz pozostałych elementów ścieżki transmisji.
Z tego względu komputer PC należy traktować przede wszystkim jako jednostkę sterującą i generującą bodźce testowe, a nie jako deterministyczne źródło sygnałów czasu rzeczywistego.

\section{Paradygmaty obiektowe i wstrzykiwanie zależności w środowisku Python}
\label{sec:python_architektura_oop}

Język Python dzięki swojej elastyczności i popularności jest często używany w szybkiej analizie danych.
Wraz z tym pojawia się wiele nieustrukturyzowanych, proceduralnych skryptów.
W przypadku autorskiego systemu HIL \textit{ESPy-Nexus}, który docelowo realizuje tysiące scenariuszy testowych, zarządza połączeniami sieciowymi, steruje portami szeregowymi i procesuje relacyjne bazy danych (SQLite), takie podejście nieuchronnie doprowadziłoby do powstania trudnego w utrzymaniu, nieskalowalnego kodu.

Aby temu zapobiec, architekturę systemu nadzorczego podzielono na odseparowane pakiety:
\begin{itemize}
    \item \texttt{core} - konfiguracja i logowanie,
    \item \texttt{runner} - silnik wykonawczy i analityczny,
    \item \texttt{control\_plane} - warstwa sterowania,
    \item \texttt{data\_plane} - warstwa transmisji danych.
\end{itemize}
Całość oparto na rygorystycznym programowaniu obiektowym (OOP) oraz wzorcach projektowych.

\subsection{Abstrakcyjne interfejsy bazowe}

Podstawą uniwersalności platformy jest narzucenie ścisłych kontraktów dla poszczególnych modułów transmisyjnych.
Wykorzystując wbudowaną bibliotekę \texttt{abc} (\textit{ang. Abstract Base Classes}), zdefiniowano abstrakcyjne interfejsy, takie jak \texttt{BaseControlPlane} oraz \texttt{BaseDataPlane}.
Zmuszają one każdą nową implementację do bezwzględnego zaimplementowania określonych metod (np. \texttt{connect()}, \texttt{disconnect()}, \texttt{transmit()}) niezależnie czy będzie to port szeregowy, czy gniazdo UDP.

Takie podejście realizuje fundamentalną zasadę otwarte-zamknięte (\textit{Open/Closed Principle}) oraz zasadę podstawienia Liskov (\textit{Liskov Substitution Principle}) z \cite{solid}.
Silnik wykonawczy wie jedynie, że operuje na obiekcie typu \texttt{BaseDataPlane}, co całkowicie uniezależnia wewnętrzną logikę testów od fizycznej implementacji poszczególnych protokołów transportowych.

\subsection{Wzorzec Fabryki i Wstrzykiwanie Zależności}
Tworzenie obiektów zarządzających warstwą fizyczną sprzętu jest procesem skomplikowanym.
Decyzja o tym, czy system użyje protokołu binarnego, tekstowego, czy wirtualnego (Mock) zależy od profilu konfiguracyjnego (\texttt{RunnerSettings}).
Logikę decyzyjną odizolowano w pliku uruchomieniowym \texttt{main.py}, stosując wzorzec \textbf{Fabryki} (\textit{Factory Pattern}).
Funkcje takie jak \texttt{build\_data\_planes()} parsują żądaną konfigurację i zwracają zainstancjowane obiekty konkretnych klas (np. \texttt{TcpDataPlane}).

Gotowe, skonfigurowane obiekty połączeniowe nie są jednak globalnie dostępne dla reszty aplikacji.
Aby utrzymać najwyższy stopień luźnego powiązania (\textit{ang. Loose Coupling}), zastosowano paradygmat \textbf{Wstrzykiwania Zależności} (\textit{ang. DI - Dependency Injection}).
Jak zaprezentowano na Listingu \ref{lst:python_di_main}, gotowe instancje warstwy sterującej i transmisyjnej przekazywane są jako argumenty do konstruktora głównej maszyny stanów (\texttt{TestEngine}).

\begin{listing}[htbp]
    \caption{Wykorzystanie wzorca Fabryki i Wstrzykiwania Zależności do izolacji logiki wykonawczej od warstwy sprzętowej (\texttt{main.py}) \cite{espy_nexus_python_host}}
    \label{lst:python_di_main}
    \begin{minted}[fontsize=\footnotesize, linenos]{python}
def main() -> None:
    # Załadowanie profilu konfiguracyjnego
    config = get_active_profile()

    # Inicjalizacja instancji sprzętowych (Wzorzec Fabryki)
    # Fabryki zwracają obiekty implementujące interfejsy BaseControlPlane i BaseDataPlane
    control_plane = build_control_plane(config)
    data_planes = build_data_planes(config)

    # Wygenerowanie wektora testowego
    test_matrix = build_matrix(config)

    # Wstrzyknięcie zależności (DI) do konstruktora silnika TestEngine
    # Silnik operuje wyłącznie na wstrzykniętych interfejsach, co pozwala 
    # na łatwą wymianę protokołów (np. na MockDataPlane) w celach testowych.
    engine = TestEngine(
        control_plane=control_plane,
        data_planes=data_planes,
        db_path=config.raw_db_path,
    )
    
    # Rozpoczęcie fazy akwizycji danych HIL
    engine.run_matrix(matrix=test_matrix)
\end{minted}
\end{listing}

Rozdzielenie etapu \textit{składania} komponentów od ich faktycznego użycia to ogromna zaleta z perspektywy inżynierii testów.
Pozwala to przełączyć system w profil deweloperski, wstrzykując obiekty typu \texttt{MockControlPlane} i \texttt{MockDataPlane}.
Silnik \texttt{TestEngine} wykona wówczas pełny proces eksperymentalny (od symulowanej wysyłki, po zapis do bazy SQLite i uruchomienie procesu analitycznego ETL), bez informacji, że po drugiej stronie nie znajduje się fizyczny mikrokontroler ESP32.
Znacząco usprawnia to rozwój oprogramowania (\textit{ang. Software Development Lifecycle}) bez konieczności angażowania środowiska sprzętowego do weryfikacji warstwy logicznej w Pythonie.
